\chapter{绪论}\label{chap:introduction}

\section{研究背景与意义}

\subsection{研究背景}

语音是人类最自然、最基本的交流方式。与文本相比，语音不仅承载了语义信息，还蕴含了说话人的情感、语气、语速和韵律等丰富的副语言信息，是人与人之间传递意图和情感的核心媒介。在全球化进程不断加速的今天，不同语言的人群之间的交流需求日益增长，语音作为面对面沟通的主要载体，跨语言的语音沟通能力对于国际合作、文化交流和信息传播具有重要意义。与此同时，随着人工智能技术的快速发展，人们对人机交互的期望也从传统的键盘文本输入，转向更为自然的语音对话形式，无论是智能手机上的语音助手、车载交互系统还是智能家居设备，语音正在成为人与机器沟通的主要界面。在这一背景下，语音到语音生成（Speech-to-Speech Generation）成为自然语言处理和语音处理领域的核心研究课题。该课题的核心科学问题在于：如何设计模型结构，使其能够以语音为输入和输出，同时兼顾生成内容的质量、语音的自然度以及响应的实时性。围绕这一问题，模型结构经历了从级联架构到端到端编码器-解码器架构，再到以大语言模型为核心的仅解码器架构的持续演进，每一次结构范式的转变都带来了新的研究机遇和技术挑战。

语音到语音生成最早采用的是级联架构，将完整的生成过程拆分为若干独立的子模块并串联执行。典型的级联系统由自动语音识别（Automatic Speech Recognition, ASR）、文本处理模块和语音合成（Text-to-Speech, TTS）三段组成。在语音翻译场景中，文本处理模块为机器翻译（Machine Translation, MT），负责将ASR输出的源语言文本翻译为目标语言文本；在语音对话场景中，文本处理模块为大语言模型（Large Language Model, LLM），负责基于ASR输出的文本进行理解和生成。无论何种场景，级联架构的结构本质是一致的：语音输入先经ASR转录为文本，文本经过处理后再由TTS合成为语音输出。这种模块化的结构设计使各子系统可以独立开发和优化，且每个子任务都有相对成熟的技术积累。目前，包括Google翻译、Apple Siri、Amazon Alexa在内的主流商用系统仍主要采用这一架构。

然而，级联架构的结构特性决定了其存在若干固有缺陷。第一，各模块之间的错误会逐级传播和累积。ASR模块的识别偏差会直接传递给下游模块，导致最终输出的质量受限于整个流水线中最薄弱的环节。第二，以文本作为模块间唯一的信息传递媒介，导致语音中蕴含的副语言信息在转录过程中大量丢失。源语音中的情感、语气、语速等特征无法通过文本传递给后续模块，最终合成的语音难以保留源说话人的风格特征。第三，多个模块的串行处理引入了较高的端到端延迟。每个模块需要等待上一个模块完成输出后才能开始工作，各环节的处理时间逐级累加，难以满足实时交互对低延迟的要求。更根本的问题在于，由于各模块使用不同的训练目标独立优化，整个系统无法进行端到端的全局优化，模块间的信息瓶颈从结构上限制了系统的性能上限。

为克服级联架构的结构局限，研究者们探索了端到端的语音到语音生成方法，旨在使用单一模型直接完成从输入语音到输出语音的映射，从结构上消除中间文本表示带来的信息瓶颈。这一阶段的模型主要基于Transformer\citep{vaswani2017attention}编码器-解码器架构构建：编码器负责将输入语音编码为连续表示序列，解码器则基于编码器的输出自回归地生成目标序列。自Translatotron\citep{translatotron}开创性地提出端到端语音翻译模型以来，研究者们在模型结构上进行了持续创新。在语音表征方面，早期模型以梅尔谱等连续声学特征作为语音的表示形式，而后续研究发现将语音通过自监督学习模型编码为离散单元\citep{lee-etal-2022-textless}可以有效降低建模难度，使语音生成问题转化为与文本生成类似的离散序列预测问题。在解码策略方面，两阶段解码方法\citep{translatotron2,inaguma-etal-2023-unity}在解码器中引入文本作为中间表示，先生成目标语言文本再合成对应语音，利用文本的高语义密度来引导语音生成，取得了较好的效果。端到端编码器-解码器架构从结构上实现了输入到输出的直接映射，理论上能够减少错误累积、保留副语言信息并实现全局优化。然而，这一架构在结构上面临两方面的核心挑战。其一，语音到语音生成任务本身极为复杂，同时涉及声学信号处理、深层语义理解和跨语言转换，单一编码器-解码器结构的建模能力有限，如何设计更有效的模型结构来兼顾内容准确性与语音自然度，是一个重要的研究问题。其二，语音序列的长度远大于文本序列，自回归逐帧解码会带来极高的计算延迟，而现有的非自回归加速方法往往以牺牲生成质量为代价，如何从结构上实现高效解码与高质量生成的兼顾，是实现实时交互的关键瓶颈。

近年来，以ChatGPT\citep{chatgpt}和GPT-4o\citep{gpt4o}为代表的大语言模型取得了突破性进展，展现出远超传统编码器-解码器模型的语言理解、推理和生成能力。这一突破推动了语音到语音生成领域的又一次结构范式转变：模型架构从编码器-解码器演变为以仅解码器（Decoder-only）大语言模型为核心的新框架。在这一框架下，大语言模型作为核心的语言处理引擎，其强大的文本能力为生成高质量、连贯的回复内容提供了基础；而语音模态的接入则通过在大语言模型的输入端和输出端分别添加语音编码器和语音解码器来实现。这种结构设计的核心优势在于能够直接复用大语言模型在大规模文本数据上习得的知识，无需从头训练完整的语音到语音生成模型。研究者们在这一框架下探索了两条技术路线：一是将语音词元直接融入大语言模型的词表，通过大规模多模态数据进行原生联合训练；二是保持大语言模型参数不变或微调，通过外挂模块以较低的训练成本扩展语音能力。无论哪条路线，仅解码器架构都面临新的结构性挑战：在输入端，如何设计语音编码器的结构使其输出能够与大语言模型的文本表示空间有效对齐；在输出端，如何设计语音解码器的结构使其既能生成高自然度的语音，又能满足实时交互对流式输出和低延迟的要求。特别是在语音解码器的设计上，非自回归结构速度快但生成质量受限，自回归结构质量高但难以实现流式输出，如何从结构上调和这一矛盾是该方向的核心难题。

仅解码器架构下的语音交互模型目前主要面向基于轮次的对话场景，即用户说完一句话后系统才开始生成回复，交互过程呈现严格的"你说我听、我说你听"的串行模式。这类模型在结构上依赖外部的语音活动检测（Voice Activity Detection, VAD）模块来判断用户发言的起止边界，而VAD仅基于声学特征进行判断，缺乏对语义信息的理解。在真实的人际对话中，双方的发言往往存在重叠、打断和自然的轮替，对话时机的把控依赖于对语义内容的深层理解。这意味着语音交互模型需要从基于轮次的半双工结构演进为全双工结构——模型在持续接收用户音频输入的同时，自主决定何时开始回复、何时停止回复。全双工建模从根本上改变了模型的序列组织方式：模型需要在每个时刻同时处理输入和输出两个语音信道，这与语言模型单向顺序生成的基本建模方式相悖。此外，为保证生成质量而引入的文本信道与语音信道之间存在粒度不匹配的问题，如何在统一的模型结构中协调多个信道的同步建模，是全双工架构设计的核心挑战。

综上，语音到语音生成领域的模型结构经历了级联架构、端到端编码器-解码器架构和以大语言模型为核心的仅解码器架构三个发展阶段，并正在向全双工架构演进。每一次结构范式的转变都解决了前一阶段的核心瓶颈，同时也引入了新的科学问题。本文从模型结构的角度出发，围绕这一演进历程中的关键技术问题展开研究：在编码器-解码器架构下，探索高效的非自回归解码策略和模块化模型组合方法；在仅解码器架构下，研究低延迟语音交互、高质量流式语音生成以及全双工建模等前沿问题，旨在设计兼顾高质量内容生成、高自然度语音合成和低响应延迟的语音到语音生成模型。

\subsection{研究意义}

本文的研究兼具理论价值和应用价值。

从理论角度看，语音到语音生成可以被视为一种深度融合内容理解与模态转换的复杂序列到序列问题，涉及两个核心科学挑战。其一是异构数据的统一建模问题：无论是编码器-解码器架构还是仅解码器架构，都需要在单一模型中有效处理离散、高语义密度的文本和连续、富含韵律信息的声学信号，如何设计模型结构来弥合两种模态之间的差异是基础性难题。其二是长序列的高效生成问题：语音信号的信息密度低、时间跨度长，如何在降低解码延迟的同时保证生成质量、实现流式输出，是序列生成领域的关键技术瓶颈。本文针对这两个挑战，在非自回归解码、模型组合、流式语音生成和全双工建模等方向上进行了系统性探索，为语音到语音生成的模型结构设计提供了新的思路。

从应用角度看，本文的研究成果可服务于两大核心场景。在跨语言交流方面，高质量、低延迟的语音翻译系统将推动国际会议同声传译、跨国商务洽谈和在线教育等领域的智能化发展，实现流畅、自然且保留情感特征的跨语言语音沟通。在人机交互方面，基于大语言模型的语音交互技术将使智能客服、车载助手等应用从简单的一问一答升级为流畅连贯的对话体验，而全双工交互能力的引入将进一步使人机对话接近人与人之间的自然交流，提升用户的参与感和信任度。


\section{国内外研究现状}

语音到语音生成领域的研究大致沿两条主线展开：基于编码器-解码器架构的端到端语音翻译，以及基于仅解码器架构的语音大语言模型。本节从这两条主线出发，简要梳理该领域的研究进展，并在此基础上分析当前研究中存在的关键问题。更为全面和深入的文献综述将在第\ref{chap:related-work}章中展开。

\subsection{端到端语音翻译}

端到端语音翻译旨在直接将源语言语音转换为目标语言的文本或语音，跳过中间的语音识别步骤。

在语音到文本翻译（Speech-to-Text Translation, S2TT）方面，核心挑战在于高质量平行数据的稀缺以及语音与文本之间的模态差异。为缓解数据瓶颈，研究者们探索了多种利用语音识别和机器翻译数据的策略，包括预训练\citep{bansal2019pre,wang-etal-2020-curriculum}、多任务学习\citep{weiss2017sequence,anastasopoulos2018tied}、知识蒸馏\citep{liu2019end,gaido2020end}和数据增强\citep{jia2019leveraging,lam-etal-2022-sample}。为弥合语音与文本之间的模态差异，研究者们提出了共享语义空间映射、对比学习\citep{ye-etal-2022-cross,ouyang-etal-2023-waco}和联合预训练\citep{tang-etal-2022-unified,Chen2022MAESTROMS}等方法，使模型能够学习跨模态的统一表示。

在语音到语音翻译（Speech-to-Speech Translation, S2ST）方面，Translatotron\citep{translatotron}开创了端到端S2ST的研究方向，直接从源语言语音生成目标语言的梅尔谱。随后，离散语音单元的引入\citep{lee-etal-2022-textless}将语音生成转化为离散序列预测问题，降低了建模难度。为提升翻译质量，两阶段解码方法\citep{translatotron2,inaguma-etal-2023-unity}先生成文本再合成语音，利用文本作为中间表示引导生成过程。在解码效率方面，非自回归方法\citep{huang2023transpeech,diffs2ut}通过并行生成加速解码，但通常伴随翻译质量的下降。如何在保持高翻译质量的同时实现高效解码，仍然是该领域的核心难题。

\subsection{语音大语言模型}

随着大语言模型展现出强大的语言理解和生成能力，如何将这些能力扩展到语音模态成为研究热点。

早期工作将语音离散化为词元序列，直接在语言模型框架下进行建模。AudioLM\citep{borsos2023audiolm}和VALL-E\citep{wang-etal-2023-valle}等模型探索了基于语音词元的语言建模方法，但这些模型未充分利用通用大语言模型的文本知识。

为更好地利用大语言模型的能力，研究者们发展出两条技术路线。第一条路线采用原生多模态训练，将语音词元直接加入大语言模型的词表，通过大规模多模态数据进行联合预训练。SpeechGPT\citep{zhang-etal-2023-speechgpt}和AnyGPT\citep{zhan2024anygpt}是这一路线的代表性工作，但这种方式需要大量的训练数据和计算资源。第二条路线采用组合式架构，通过外挂语音编码器赋予大语言模型语音理解能力，再通过语音解码器实现语音输出。Qwen-Audio\citep{Qwen-Audio}、SALMONN\citep{tang2024salmonn}等模型在语音理解任务上取得了优异表现。在完整的语音交互方面，GLM-4-Voice\citep{glm4voice}通过语音与文本的交织训练实现了端到端的语音对话，Mini-Omni\citep{xie2024miniomni}和Moshi\citep{kyutai2024moshi}则探索了同时生成文本和语音的方法。

在全双工交互方面，传统的基于轮次的模型依赖外部语音活动检测来切分用户发言，无法支持打断、停顿等灵活的交互模式。Moshi\citep{kyutai2024moshi}通过并行建模用户和模型两个语音流迈出了全双工交互的第一步，但其将两个语音流简单并联的方式限制了模型对对话时机的精细控制。如何使模型在持续接收音频的同时自主决策回复时机，实现自然的全双工对话，仍是一个开放的研究问题。

\subsection{研究现状分析}

综合以上研究现状，可以发现当前语音到语音生成领域存在以下五个尚未解决的关键问题。

第一，在端到端语音翻译中，自回归模型的翻译质量较高但解码速度慢，现有的非自回归模型虽然解码快但翻译质量显著下降，如何在非自回归框架下实现与自回归模型相当的翻译质量，是提升语音翻译效率的关键。

第二，端到端语音翻译模型需要从头训练，难以复用语音到文本翻译和语音合成领域已有的高质量预训练模型，其根本原因在于不同模型之间的词表和表示空间不兼容，如何打通这一接口壁垒以灵活组合现有模型，是降低训练成本和提升翻译性能的重要途径。

第三，将大语言模型扩展为语音交互系统时，如何以较低的训练成本同时赋予模型语音输入和语音输出能力，并实现低延迟的实时响应，是构建实用语音交互系统面临的核心挑战。

第四，在组合式语音大语言模型中，非自回归语音解码器速度快但生成语音的自然度不足，自回归语音解码器质量高但难以实现流式输出，如何在保证语音自然度的同时实现实时的流式语音生成，是提升语音交互体验的瓶颈问题。

第五，当前语音交互模型均面向基于轮次的对话场景，依赖外部的语音活动检测模块来判断用户发言边界。在真实对话中，双方的发言存在重叠和灵活交替，对话时机的把控需要语义层面的深层理解。如何使模型在持续接收音频的同时自主决定回复时机，实现全双工语音交互，是迈向自然人机对话的关键挑战。


\section{研究内容}

针对上述五个关键科学问题，本文开展了五项研究工作，从编码器-解码器架构下的语音翻译逐步推进到仅解码器架构下的语音交互和全双工对话。各项工作之间存在紧密的技术递进关系：前两项工作在编码器-解码器框架下分别解决了非自回归解码和模型组合的问题；后三项工作在大语言模型框架下依次解决了低延迟语音交互、高质量流式语音生成和全双工建模的问题。五项研究内容如下。

\subsection{基于有向无环图的非自回归语音到语音翻译}

针对第一个科学问题，本文第\ref{chap:daspeech}章提出DASpeech，一种基于有向无环图（Directed Acyclic Graph, DAG）Transformer的非自回归语音到语音翻译方法。现有的非自回归S2ST模型在解码速度上虽有优势，但翻译质量与自回归模型存在较大差距。DASpeech采用两阶段解码框架，第一阶段使用DA-Transformer作为语言解码器，以有向无环图结构建模目标序列的多种可能路径，通过前向-后向算法在所有路径上计算期望隐状态并传递给第二阶段的非自回归声学解码器FastSpeech 2进行语音合成。实验结果表明，DASpeech在翻译质量上达到甚至超过了自回归模型的水平，同时实现了高达18.53倍的解码加速。相关成果已发表于NeurIPS 2023\citep{fang2023daspeech}。

\subsection{基于词表转换的可组合语音到语音翻译}

针对第二个科学问题，本文第\ref{chap:comspeech}章提出ComSpeech，一种通过词表转换实现模型组合的语音到语音翻译框架。现有的端到端S2ST模型需要从头训练，难以利用语音到文本翻译和语音合成领域的高质量预训练模型，其根本障碍在于不同模型使用不同的词表（如子词词表与音素词表），输出空间不兼容。ComSpeech提出基于连接时序分类（Connectionist Temporal Classification, CTC）的词表转换器，以非自回归方式将S2TT模型的输出序列转换为TTS模型的输入序列，从而将任意S2TT和TTS模型桥接为完整的S2ST系统。此外，ComSpeech还探索了零样本训练模式，无需任何平行语音翻译数据即可构建S2ST系统。实验结果表明，ComSpeech显著超越了现有的端到端基线模型，零样本模式也取得了接近有监督训练的性能。相关成果已发表于ACL 2025\citep{fang2024comspeech}。

\subsection{基于非自回归语音解码的实时语音交互大模型}

针对第三个科学问题，本文第\ref{chap:llama-omni}章提出LLaMA-Omni，一种基于大语言模型的低延迟语音交互模型。现有方法将语音交互分解为语音识别、语言模型推理和语音合成三个串行阶段，累积延迟较高。LLaMA-Omni将语音编码器、大语言模型和非自回归语音解码器集成为统一模型，其中语音解码器基于CTC在大语言模型生成文本的同时并行预测离散语音单元，实现了文本和语音的同步生成。为训练该模型，本文构建了包含200000条语音指令数据的InstructS2S-200K数据集。实验结果表明，LLaMA-Omni在内容质量和语音自然度上均优于同期基线模型，响应延迟低至226毫秒。相关成果已发表于ICLR 2025\citep{fang2024llama}。

\subsection{基于自回归流式语音解码的实时语音交互大模型}

针对第四个科学问题，本文第\ref{chap:llama-omni2}章提出LLaMA-Omni 2，在LLaMA-Omni的基础上引入自回归流式语音解码器以提升语音生成质量。LLaMA-Omni的非自回归语音解码器虽然速度快，但生成语音的准确性和自然度有限。LLaMA-Omni 2采用独立的自回归语音语言模型作为语音解码器，通过"读R写W"的流式策略实现大语言模型文本生成与语音合成的交替进行：大语言模型每生成R个文本词元后，语音解码器即生成W个对应的语音词元。为进一步提升语音质量，模型使用因果流匹配模型和HiFi-GAN声码器将离散语音词元转换为连续语音波形。实验结果表明，LLaMA-Omni 2在语音准确性和自然度上取得了显著提升，同时保持了实时交互所需的低延迟。相关成果已发表于ACL 2025\citep{fang2025llamaomni}。

\subsection{基于全双工建模的自然语音交互大模型}

针对第五个科学问题，本文第\ref{chap:full-duplex}章提出一种基于多信道交织的全双工语音交互模型。前述语音交互模型均面向基于轮次的对话场景，依赖外部VAD模块进行发言边界检测，无法支持打断、停顿等灵活的交互模式。本章将用户语音、模型回复文本和模型回复语音三个信道按固定比例组织成交织序列，统一由大语言模型进行自回归建模。文本信道中引入状态标记来编码对话状态，使全双工场景下的时机决策被自然转化为下一词元预测问题，无需额外的分类模块。在训练策略上，基于已有的语音对话模型进行轻量的全双工微调，并引入直接偏好优化进一步强化模型的时机决策能力。实验结果表明，该模型在轮替和打断场景上均取得了较高的成功率，对话质量相比基线模型取得了显著提升。


\section{章节安排}

本文共分为八章，各章内容安排如下：

第一章为绪论，介绍本文的研究背景与意义、国内外研究现状、关键科学问题和研究内容。

第二章为研究现状，全面梳理语音到语音生成领域的相关工作，涵盖语音表征、建模范式、语音到文本翻译、语音到语音翻译和语音大语言模型等方面的研究进展，并通过综合分析指出当前存在的五个关键科学问题。

第三章介绍本文的第一项工作DASpeech。针对非自回归语音翻译中翻译质量不足的问题，提出基于有向无环图Transformer的两阶段非自回归语音到语音翻译方法，在翻译质量和解码速度之间实现了良好的平衡。

第四章介绍本文的第二项工作ComSpeech。针对端到端语音翻译模型难以复用预训练模型的问题，提出基于CTC词表转换器的可组合语音翻译框架，实现了对任意语音到文本翻译模型和语音合成模型的灵活组合。

第五章介绍本文的第三项工作LLaMA-Omni。针对语音交互系统延迟高的问题，提出基于非自回归CTC语音解码器的语音交互大模型，实现了文本与语音的同步生成和低至226毫秒的响应延迟。

第六章介绍本文的第四项工作LLaMA-Omni 2。针对非自回归语音解码器生成质量受限的问题，提出基于自回归流式语音解码器的语音交互大模型，在保持实时交互的同时大幅提升了语音的准确性和自然度。

第七章介绍本文的第五项工作，基于全双工建模的自然语音交互大模型。针对基于轮次的模型无法支持灵活交互的问题，提出基于多信道交织的全双工语音交互模型，使模型能够在持续接收音频输入的同时自主决策回复时机。

第八章为总结与展望，总结本文的主要研究工作和贡献，并对未来的研究方向进行展望。
